Fix \(m\geq1\) and rational \(0<\epsilon<1/2\), and write \(N=N_m\) and \(K=K_{\epsilon}\).  For \(q\in K\), put
\[
 d(q)=p(q)\{1-p(q)\},\qquad
 a(q)=q_{11}\{1-p(q)\}-q_{01}p(q).
\tag{1}
\]
Then \(d(q)\geq\epsilon(1-\epsilon)>0\) and \(f(q)=a(q)/d(q)\).

For \(1\leq k\leq N\), let \(\Phi_k(b,t)\) be the assertion that there exist \(q^1,\ldots,q^k\in K\) and \(\omega_1,\ldots,\omega_k\geq0\) satisfying
\[
 \sum_{j=1}^k\omega_j=1,\qquad
 b_c=\sum_{j=1}^k\omega_jB_c^{(m)}(q^j)\quad(c\in\mathcal I_m),
\tag{2}
\]
and
\[
 t\prod_{j=1}^kd(q^j)
 =\sum_{j=1}^k\omega_ja(q^j)
   \prod_{\substack{1\leq r\leq k\\r\neq j}}d(q^r).
\tag{3}
\]
The strict positivity in (1) makes (3) equivalent to \(t=\sum_j\omega_jf(q^j)\).  Thus \(\Phi_k\) is a first-order formula involving only polynomial equalities and weak inequalities with rational coefficients.  Zero weights and repeated locations show that it expresses existence of an atomic witness with support size at most \(k\).

Let \(\Phi=\Phi_N\).  Every \(\Phi(b,t)\) witness belongs to \(\mathcal F_{m,\epsilon}(b)\).  Conversely, the all-feasible endpoint result in the sharp primal--dual theorem gives, for every \(b\in\mathcal B_{m,\epsilon}\), measures with at most \(N\) atoms attaining both \(L_m(b)\) and \(U_m(b)\).  Hence
\[
 \Psi(b):\Longleftrightarrow\exists t\,\Phi(b,t)
 \quad\Longleftrightarrow\quad b\in\mathcal B_{m,\epsilon},
\tag{4}
\]
and
\[
 \min\{t:\Phi(b,t)\}=L_m(b),\qquad
 \max\{t:\Phi(b,t)\}=U_m(b).
\tag{5}
\]
The set parameterized by \(\Phi\) is compact, because it is the continuous image of the compact space \(\Delta_{N-1}\times K^N\).  Therefore the minima and maxima in (5) are attained.  The CAD quantifier-elimination lemma applied to (4) gives a quantifier-free semialgebraic definition of the moment body.

We next define its relative boundary without confusing it with the ambient boundary.  A first-order definition of \(A=\operatorname{aff}(\mathcal B_{m,\epsilon})\subseteq\mathbb R^N\) is
\[
 \operatorname{Aff}(x):\Longleftrightarrow
 \exists z^0,\ldots,z^N\ \exists\alpha_0,\ldots,\alpha_N
 \left[
 \bigwedge_{j=0}^N\Psi(z^j),\quad
 \sum_{j=0}^N\alpha_j=1,\quad
 x=\sum_{j=0}^N\alpha_jz^j
 \right].
\tag{6}
\]
The right side is contained in the affine hull.  Conversely, any affine representation in \(\mathbb R^N\) using more than \(N+1\) nonzero coefficients can be shortened by affine dependence while preserving both the coefficient sum and the represented point.  Thus (6) is exactly \(A\).  Define
\[
 \begin{split}
 \operatorname{RI}(b):\Longleftrightarrow{}&\Psi(b)\ \wedge\
 \exists\rho>0\ \forall x\left[
 \operatorname{Aff}(x)\ \wedge\ \sum_c(x_c-b_c)^2<\rho^2
 \Rightarrow\Psi(x)
 \right],\\
 \operatorname{RB}(b):\Longleftrightarrow{}&\Psi(b)\ \wedge\neg\operatorname{RI}(b).
 \end{split}
\tag{7}
\]
Formula \(\operatorname{RI}\) is precisely relative interior in \(A\).  The body is compact, since (4) makes it the projection of a compact parameterized set, and is therefore closed in \(A\).  Consequently \(\operatorname{RB}\) defines exactly \(\operatorname{relbd}(\mathcal B_{m,\epsilon})\).  Quantifier elimination makes all three predicates in (6)--(7) semialgebraic.

Define the endpoint graphs by
\[
 \begin{split}
 G_L(b,t)&:\Longleftrightarrow
 \Phi(b,t)\ \wedge\ \forall s\{\Phi(b,s)\Rightarrow t\leq s\},\\
 G_U(b,t)&:\Longleftrightarrow
 \Phi(b,t)\ \wedge\ \forall s\{\Phi(b,s)\Rightarrow s\leq t\}.
 \end{split}
\tag{8}
\]
By compactness and (5), these graphs are nonempty over every feasible \(b\), and their unique values are \(L_m(b)\) and \(U_m(b)\).  Therefore
\[
 P_0(b):\Longleftrightarrow\exists t\{G_L(b,t)\wedge G_U(b,t)\}
\tag{9}
\]
is exactly the point-identification predicate \(L_m(b)=U_m(b)\).  Since the full sharp fiber is the nonempty interval \([L_m(b),U_m(b)]\), its affine dimension is zero when (9) holds and one otherwise.

For \(e\in\{L,U\}\), \(1\leq k\leq N\), define
\[
 E_{e,k}(b):\Longleftrightarrow
 \exists t\{G_e(b,t)\wedge\Phi_k(b,t)\},
 \qquad E_{e,0}(b):\Longleftrightarrow\mathrm{false}.
\tag{10}
\]
These are semialgebraic after quantifier elimination.  The endpoint compression theorem gives \(E_{e,N}(b)\) for every feasible \(b\).  Hence the minimum endpoint support size is \(k\) exactly when
\[
 S_{e,k}(b):\Longleftrightarrow
 E_{e,k}(b)\wedge\neg E_{e,k-1}(b),
 \qquad 1\leq k\leq N.
\tag{11}
\]
The negated term in (11) certifies that no smaller support can attain that endpoint.

Complete \(\mathcal A_{m,\epsilon}\) by adjoining the formulas (7), (9), and (11) to its face stratification, eliminating all quantified variables, and taking a common truth-invariant CAD in the \(b\)-space.  Retain exactly the cells satisfying \(\operatorname{RB}\).  There are finitely many retained semialgebraic cells.  On each, (9), exactly one \(S_{L,k}\), and exactly one \(S_{U,l}\) have constant truth values.  Thus point identification, fiber dimension zero or one, and the ordered pair of minimum endpoint support sizes are constant on each cell, with both sizes in \(\{1,\ldots,N_m\}\).

All coefficients are rational because \(m\) is fixed and \(\epsilon\) is rational.  The CAD lemma proves termination and produces defining sign formulas for every nonempty retained cell.  Its witness extraction applied to (8) and (10) returns endpoint values, atom locations, and weights.  If a returned \(k\)-slot representation had a zero weight or two equal positive-weight locations that could be consolidated, it would give a \((k-1)\)-slot representation, contradicting \(S_{e,k}\); hence a witness on that stratum has the certified minimum support size.  For any separately supplied algebraic point of a stratum, specialization of the same formulas and a further terminating CAD run returns its endpoint witnesses.  For rational boundary inputs this agrees with the boundary-capable branch of the exact CAD recovery theorem.

Finally, the boundary Dirac-regression proposition verifies that the point-identified alternative is genuinely needed.  For \(m\geq2\), it gives
\[
 \mathcal F_{m,\epsilon}(g_m(q^*))=\{\delta_{q^*}\},
 \qquad L_m(g_m(q^*))=U_m(g_m(q^*))=f(q^*).
\tag{12}
\]
Such a moment vector is extreme: a nontrivial convex decomposition would, after mixing representing measures for its two terms, contradict the singleton statement in (12).  Moreover, differentiation of the multinomial identity gives
\[
 q_{ay}=\frac1m\sum_{c\in\mathcal I_m}c_{ay}B_c^{(m)}(q),
\tag{13}
\]
so \(g_m\) is injective.  The overlap slab contains distinct perturbations of the uniform vector that keep propensity \(1/2\), so the moment body is not a singleton.  An extreme point of a non-singleton convex set is not in its relative interior.  Thus (12) lies on the relative boundary and is assigned the zero-dimensional alternative by (9).  This completes the claimed finite boundary classification and certified witness construction.